\subsection{An exact coefficient certificate for the pure degree-eleven direction}
\label{subsec:fr-pure-eleven-certificate}

The next normalization uses a specific integral vector in the
degree-eleven rational leading space. This subsection gives all
finite data needed to identify it; it does not rely on assuming
that an arbitrary infinitesimal solution lifts to profinite GT.

Use the lexicographic Lyndon basis, with $S<T$ and the longest
proper Lyndon suffix bracketing. Let $N=(N_0,\ldots,N_{185})$ be
the integer vector whose entries are listed below, row by row,
and let $L_{11}$ be the corresponding Lie polynomial of degree
eleven:
\[
\begin{array}{rrrrrr}
8448 & -42240 & 152064 & 126720 & -236544 & 183632 \\
-635824 & -253440 & 147840 & -337248 & -37328 & -1013760 \\
1632032 & -259888 & 1144704 & 354816 & -25344 & 132496 \\
-126672 & 2133568 & -1739824 & 1611760 & 7131552 & 1700885 \\
31568 & -2423936 & 3414176 & 1545536 & -2644688 & 195408 \\
-1665643 & -354816 & -59136 & -713392 & 194976 & 145696 \\
4456096 & -2655024 & -2873338 & 428624 & -1785296 & -2847760 \\
-1650127 & 283648 & -2631426 & -10156516 & -2055701 & -246992 \\
3755168 & -813519 & 5663888 & -8344398 & -11894626 & -2047685 \\
-116048 & 2904305 & -3569317 & 2634032 & 1206347 & 2650091 \\
1319275 & 253440 & -53360 & 318784 & 782288 & -238768 \\
-4135792 & -2459504 & -4508616 & -267936 & -883885 & 6968697 \\
973418 & -163168 & -313456 & 5564885 & -16000 & -2331033 \\
6357753 & -2439271 & 4203469 & 16940303 & 673697 & -188229 \\
6915145 & 3851869 & 1774080 & -411952 & -1684752 & -457131 \\
-5215638 & 71769 & -6124789 & 11170493 & 1866801 & 1075125 \\
4395714 & 6237830 & 3319616 & -2154437 & 718591 & -2801034 \\
6313939 & 13062190 & -859952 & 943076 & 1130196 & -7260032 \\
42853 & -1695184 & -375936 & -126720 & -103312 & 4819387 \\
2892779 & 11269886 & -130547 & -2565834 & 2461246 & -120507 \\
1218785 & -2554734 & 137392 & 3648115 & -2401133 & -1955433 \\
-1373195 & -1404391 & 2248757 & -7660926 & -1349408 & 2951056 \\
1650701 & 1642820 & 1621284 & 1329456 & 1895999 & -4035212 \\
-7424912 & -6067296 & 130704 & -817232 & 1755411 & 1346142 \\
1712064 & -4708057 & -5711589 & -6562098 & 5671072 & -7032843 \\
5243088 & 2583680 & 1328720 & -2287408 & 1452176 & -3228615 \\
1958704 & -2843056 & 1873664 & -2198032 & -3771728 & -51344 \\
1149616 & -372592 & 1042688 & -163840 & 10672 & 42240 \\
2783575 & -402752 & -1588080 & 1981456 & -300592 & 6888560 \\
-2296000 & -830880 & 1240240 & 3721840 & -699120 & 143616 \\
1631792 & 950752 & 321024 & 375936 & 181632 & -8448 \\
\end{array}
\]
In particular $N_0=8448$. Let
$M=(M_0,\ldots,M_{185})$ be the integral Lyndon vector of
$\{P_3,L_8\}$, already specified by the finite bracket recursion.
Its depth-one coefficient is zero and its $b_{11}$ coefficient
is eleven.

\begin{lemma}\label{lem:fr-pure-eleven-certificate}
For every full-GT element whose dyadic word starts in degree
eleven, with degree-eleven leading term $q$, one has
\[
 q=\frac{[S^{10}T]q}{8448}L_{11}
       +\frac{[S^6TST^3]q}{11}\{P_3,L_8\}.
\]
In particular $[S^{10}T]L_{11}=8448$,
$[S^6TST^3]L_{11}=0$, and $L_{11}$ is primitive as an
integral $2$-adic coefficient vector.
\end{lemma}

\begin{proof}
Here is an entirely finite necessary-equation certificate.
In logarithmic variables write
$\log f=q+r+O(13)$, with $q$ homogeneous of degree eleven
and $r$ homogeneous of degree twelve. The associated Lyndon
coordinates form $521=186+335$ columns, ordered first by degree
and then lexicographically, with indices $0,\ldots,520$.
Define
\[
 \begin{split}
 I(a)&=a(S,T)+a(T,S),\\
 C(a)&=a(S,T)+a(-S-T,S)+a(T,-S-T),\\
 J(a)&=[\varepsilon]\bigl(a(-S-T-\varepsilon[S,T],S)
                 +a(T,-S-T-\varepsilon[S,T])\bigr).
 \end{split}
\]
The full duality and multiplier-one cycle imply
\[
 I(q)=0,\qquad I(r)=0,\qquad C(q)=0,\qquad
                         J(q)+2C(r)=0.
\]
Indeed $\log((xy)^{-1})=-S-T-\tfrac12[S,T]+O(3)$,
and products of two nonconstant terms of $f$ start in degree
twenty-two. Thus the last equation is twice the degree-twelve
cycle coefficient, with no quadratic-in-$f$ contribution.

Form the integer matrix $A_0$ by extracting every nonzero
associative word coefficient of these four expressions, in
the displayed order, with words lexicographic within each
degree. This gives $12092$ rows and $521$ columns.
To these append seven necessary leading-pentagon rows.
For free Lie generators $X_1,\ldots,X_4$, define
\[
 \delta_{ij}(X_i)=[X_j,X_i],\qquad
 \delta_{ij}(X_j)=[X_i,X_j],\qquad
 \delta_{ij}(X_k)=0\quad(k\notin\{i,j\}).
\]
These derivations satisfy the infinitesimal braid relations.
For a Lie polynomial $a$ write
\[
\begin{split}
\mathcal P(a)={}&a(\delta_{12},\delta_{23}+\delta_{24})
 +a(\delta_{13}+\delta_{23},\delta_{34})
 -a(\delta_{23},\delta_{34})\\
 &-a(\delta_{12}+\delta_{13},\delta_{24}+\delta_{34})
 -a(\delta_{12},\delta_{23}).
\end{split}
\]
Take the coefficients of the following words in
$\mathcal P(q)(X_1)$, where a digit denotes its corresponding
free generator:
\[
\begin{gathered}
111111123244,\quad111112323244,\quad111123232344,\quad
111123234234,\\
111232342344,\quad112323423244,\quad121232323444.
\end{gathered}
\]
They have zero entries in the degree-twelve columns. These are
genuine necessary equations because the leading term of a
full-GT word satisfies the infinitesimal pentagon; the
derivation representation only evaluates that equation.
This fully defines a $12099$ by $521$ integer matrix $A$.

The following minor and identity are an exact integer/rational
certificate. Put
\[
 C=\{0,1,\ldots,408\}\setminus\{180,185\}.
\]
The row-index set $R$, using zero-based indices of the matrix
just defined, is the following list read row by row:
{\scriptsize
\[
\begin{array}{rrrrrrrrrrrr}
0 & 2 & 4 & 6 & 8 & 10 & 12 & 14 & 16 & 18 & 20 & 22 \\
24 & 26 & 28 & 30 & 32 & 34 & 36 & 38 & 40 & 42 & 44 & 46 \\
48 & 50 & 52 & 54 & 56 & 58 & 60 & 66 & 68 & 70 & 72 & 74 \\
76 & 78 & 80 & 82 & 84 & 86 & 88 & 90 & 92 & 98 & 100 & 102 \\
104 & 106 & 108 & 112 & 114 & 116 & 118 & 120 & 136 & 138 & 140 & 146 \\
148 & 150 & 152 & 154 & 156 & 162 & 164 & 166 & 168 & 170 & 172 & 178 \\
180 & 184 & 198 & 200 & 202 & 204 & 210 & 212 & 216 & 218 & 228 & 232 \\
292 & 298 & 300 & 306 & 308 & 330 & 332 & 338 & 340 & 2046 & 2048 & 2050 \\
2052 & 2054 & 2056 & 2058 & 2060 & 2062 & 2064 & 2066 & 2068 & 2070 & 2072 & 2074 \\
2076 & 2078 & 2080 & 2082 & 2084 & 2086 & 2088 & 2090 & 2092 & 2094 & 2096 & 2098 \\
2100 & 2102 & 2104 & 2106 & 2111 & 2113 & 2115 & 2117 & 2119 & 2121 & 2123 & 2125 \\
2127 & 2129 & 2131 & 2133 & 2135 & 2137 & 2139 & 2141 & 2143 & 2145 & 2147 & 2149 \\
2151 & 2153 & 2155 & 2157 & 2159 & 2161 & 2163 & 2165 & 2167 & 2175 & 2177 & 2179 \\
2181 & 2183 & 2185 & 2186 & 2188 & 2190 & 2192 & 2194 & 2196 & 2198 & 2204 & 2206 \\
2208 & 2210 & 2212 & 2214 & 2217 & 2219 & 2221 & 2223 & 2225 & 2227 & 2229 & 2237 \\
2239 & 2241 & 2243 & 2245 & 2248 & 2250 & 2252 & 2254 & 2256 & 2258 & 2260 & 2266 \\
2268 & 2270 & 2272 & 2274 & 2276 & 2279 & 2281 & 2312 & 2314 & 2316 & 2318 & 2320 \\
2322 & 2328 & 2330 & 2332 & 2334 & 2336 & 2338 & 2344 & 2346 & 2349 & 2351 & 2363 \\
2365 & 2367 & 2369 & 2375 & 2377 & 2380 & 2382 & 2384 & 2390 & 2392 & 2394 & 2395 \\
2397 & 2399 & 2405 & 2431 & 2437 & 2439 & 2442 & 2444 & 2446 & 2454 & 2457 & 2459 \\
2461 & 2467 & 2469 & 2488 & 2490 & 2500 & 2615 & 2617 & 2623 & 2625 & 2639 & 2645 \\
2647 & 2686 & 2709 & 2715 & 2748 & 2841 & 6016 & 6018 & 6022 & 6028 & 6030 & 6036 \\
6040 & 6046 & 6052 & 6056 & 6058 & 6064 & 6066 & 6084 & 6088 & 6090 & 6096 & 6098 \\
6102 & 6114 & 6118 & 6152 & 6154 & 6160 & 6162 & 6166 & 6176 & 6178 & 6182 & 6214 \\
6306 & 8060 & 8062 & 8064 & 8066 & 8068 & 8070 & 8072 & 8076 & 8078 & 8080 & 8082 \\
8084 & 8086 & 8088 & 8092 & 8094 & 8096 & 8098 & 8100 & 8102 & 8104 & 8106 & 8108 \\
8110 & 8112 & 8114 & 8116 & 8118 & 8126 & 8128 & 8130 & 8132 & 8134 & 8136 & 8138 \\
8140 & 8142 & 8144 & 8146 & 8148 & 8150 & 8152 & 8155 & 8157 & 8159 & 8161 & 8163 \\
8165 & 8167 & 8169 & 8171 & 8175 & 8191 & 8195 & 8197 & 8199 & 8202 & 8204 & 8206 \\
8208 & 8210 & 8212 & 8214 & 8220 & 8222 & 8226 & 8228 & 8230 & 8234 & 8236 & 8238 \\
8240 & 8242 & 8244 & 8246 & 8254 & 8256 & 8258 & 8260 & 8262 & 8266 & 8268 & 8330 \\
8332 & 8334 & 8336 & 8338 & 8346 & 8348 & 8352 & 8354 & 8356 & 8364 & 8384 & 8386 \\
8388 & 8396 & 8400 & 8402 & 8404 & 8410 & 8412 & 8415 & 8417 & 8419 & 8451 & 8638 \\
8648 & 8710 & 8734 & 8774 & 12092 & 12093 & 12094 & 12095 & 12096 & 12097 & 12098 &  \\
\end{array}
\]}
It has $407$ entries. For $B=A_{R,C}$ the determinant is
\[
 \begin{aligned}
 \det B={}&570662177953320090405597978960021340179106\cdot10^{36}\\
 &+875700910966765341012729477988352000.
 \end{aligned}
\]
The $b_{11}$ functional on the degree-eleven Lyndon coordinates
is
\[
 \ell=e_{11}-4e_7,
\]
extended by zero on the degree-twelve columns. Define the
$186$ by $521$ integer matrix $T$ by its rows
\[
 T_j=92928e_j-11N_j e_0-8448M_j\ell
       \qquad(0\le j<186).
\]
Direct integer matrix multiplication gives the certificate
\begin{equation}\label{eq:fr-pure-eleven-row-certificate}
 T_{*,C}\operatorname{adj}(B)A_{R,*}=(\det B)\,T.
\end{equation}
All entries here are determined by the printed vectors,
bracket recursions, word lists, and row/column sets. Thus
\eqref{eq:fr-pure-eleven-row-certificate} is a finite exact
coefficient identity, with no floating-point or modular
inference.

For the coefficient vector $(q,r)$ arising from a full-GT element, one has
$A(q,r)=0$. Since $\det B\ne0$, the certificate gives
$T(q,r)=0$. Dividing its $j$th equality by $92928=8448\cdot11$
gives precisely the stated leading-term formula.
The displayed vector has $N_{11}-4N_7=0$ and includes odd
entries, so its two claimed normalizations and primitivity
also follow directly.
\end{proof}
